\chapter*{Závěr}
\addcontentsline{toc}{chapter}{Závěr}

Cílem této bakalářské práce bylo navrhnout a implementovat herní prostředí pro hru inspirovanou deskovou hrou Logix a vytvořit 
ukázkové agenty, kteří využívají různé algoritmy pro rozhodování. Důraz byl kladen na implementaci a trénování neuronové sítě
pro účely ohodnocování herních pozic a její následné využití v algoritmech prohledávání herního stromu.

V rámci této práce byla vytvořena implementace v prostředí Unity, která umožňuje partii dvou lidských hráčů, hru proti počítačovým agentům
i možnost ukázkové partie dvou agentů. Současně bylo v jazyce Python vytvořeno prostředí určené pro experimenty, implementaci agentů 
a trénování neuronové sítě. 

Implementovaní agenti byli založeni na odlišných principech rozhodování. Kromě základních agentů s náhodným a heuristickým rozhodováním byli implementováni 
také agenti využívající algoritmus minimax s alfa-beta ořezáváním a algoritmus Monte Carlo Tree Search. Oba pokročilejší agenti sdílejí
společnou neuronovou síť, která poskytuje odhad kvality herních pozic a také apriorní pravděpodobnosti jednotlivých tahů.

Součástí práce byla také implementace trénovací pipeline založené na metodě self-play. Přestože se neuronovou síť podařilo natrénovat
a do obou algoritmů integrovat, experimenty ukázaly, že z hlediska kvality nedosahuje úrovně navržené heuristiky. 
Heuristický agent dosáhl nejlepších výsledků, zatímco agenti s algoritmy prohledávání herního stromu s neuronovou sítí představují
pouze funkční, ale stále nedostatečné řešení. Výsledky naznačují, že zvolená architektura neuronové sítě nebo způsob jejího trénování nebyly pro 
danou úlohu dostatečně vhodné, jelikož se s časem trénování její kvalita nezlepšila.

Možností pro další rozvoj je mnoho, zejména ale v oblasti neuronových sítí. Nabízí se delší trénování, předělání architektury či lepší 
a rychlejší implementace algoritmu MCTS pro rychlejší vyhledávání a tím i rychlejší trénování.

Přestože se nepodařilo vytvořit agenta, který by překonal heuristického agenta, byly splněny hlavní cíle práce. Vzniklo funkční
prostředí, sada vzorových agentů i infrastruktura pro trénování neuronové ohodnocovací funkce, což představuje vhodný základ pro případný
výzkum a budoucí rozvoj této hry.